\section{エージェントの「制御面」を巡る議論}
複数の実務家投稿では、コーディングエージェントを単なる補完機能ではなく、仕事全体を遂行する主体として測定・監督する方向が強調された。複数エージェントのcontrol room、open training recipe、ローカル実行可能な高性能モデル、エージェント向けウォレットや小額決済などが関連テーマとして並ぶ。単一発表ではないものの、能力評価から運用制御へ関心が移っている傾向として記録された。
\dailyxnote{単一の一次発表ではなく、複数投稿を束ねたトレンド項目。}
